\section{Front-end}

El desarrollo front-end se refiere al proceso de diseño, estructuración e implementación de las interfaces visuales con las que los usuarios interactúan dentro de un sistema. Su propósito principal es facilitar la interacción entre el usuario y las funcionalidades del sistema, permitiendo ejecutar acciones previamente definidas de manera intuitiva y eficiente \parencite{FrontentGG}.

Este proceso implica considerar diversos factores, tales como los dispositivos en los que se visualizará el sistema, así como los entornos de ejecución, en el caso de las aplicaciones web (WebApps). Asimismo, el desarrollo front-end promueve la creación de componentes reutilizables, lo que contribuye a evitar la duplicidad de código y optimiza tanto el mantenimiento como la escalabilidad del sistema.


% =======================
% |   LENGUAJES WEB     |
% =======================
\subsection{Lenguajes de desarrollo web}

Los lenguajes de desarrollo web son aquellos diseñados específicamente para la creación de sistemas que operan en entornos web, los cuales pueden ser procesados e interpretados por los navegadores o servidores. Estos lenguajes permiten construir aplicaciones accesibles a través de internet, integrando tanto la presentación como la lógica del sistema.

Dentro del desarrollo web se distinguen dos componentes principales: el front-end, que corresponde a la parte del sistema que se ejecuta en el navegador y con la que el usuario interactúa directamente, y el back-end, que se ejecuta en el servidor y se encarga de la lógica de negocio, el procesamiento de datos y la comunicación con bases de datos. Entre los lenguajes fundamentales del desarrollo web se encuentran HTML, CSS y JavaScript \ref{fig:LenguajesWeb}.

\begin{figure}[h]
  \centering
  
\includegraphics[width=0.6\linewidth]{lenguajesWeb.png}
  \caption{Lenguajes de desarrollo web.}
  \label{fig:LenguajesWeb}
\end{figure}


% =======================
% |        HTML         |
% =======================
\subsubsection{HTML (HyperText Markup Language)}

Es el lenguaje de marcado estándar utilizado para la estructuración y maquetación de páginas web. Permite definir la organización del contenido que será interpretado por los navegadores, constituyendo la base sobre la cual se construyen las interfaces de usuario en aplicaciones web \parencite{HTMLmdn}.

Se hace uso de etiquetas (tags) (Ver ejemplo \ref{lst:ejeHTML}), las cuales delimitan los distintos elementos de la página. En su mayoría, estas etiquetas están compuestas por una etiqueta de apertura y una de cierre, aunque existen algunas etiquetas autocontenidas. Cada elemento puede incluir atributos que proporcionan información adicional o configuran su comportamiento.

Estos atributos permiten la integración y manipulación de los elementos HTML a través de otros lenguajes complementarios como CSS (Cascading Style Sheets), encargado de la presentación visual, y JavaScript, responsable de la interactividad y la lógica del lado del cliente.

\begin{lstlisting}[language=HTML, caption={Ejemplo de código HTML.}, label={lst:ejeHTML}]
<html lang="en">
<head>
    <meta charset="UTF-8">
    <meta name="viewport" content="width=device-width, initial-scale=1.0">
    <title>Pagina web</title>
</head>
<body>
   <div>
    <img src="img.jpg" alt="imagen de prueba" />
   </div> 
</body>
</html>
\end{lstlisting}

% =======================
% |         CSS         |
% =======================

\subsubsection{CSS (Cascading Style Sheets)}

Es el lenguaje utilizado para definir la presentación visual de los elementos estructurados mediante HTML en una página web. Permite aplicar estilos como colores, tipografías, tamaños, márgenes, alineaciones y posicionamiento, con el objetivo de mejorar la apariencia y la experiencia de usuario en la interfaz \parencite{CSSmdn}.

La aplicación de estilos en CSS se realiza mediante selectores, los cuales permiten identificar elementos HTML a través de nombres de etiqueta, clases, identificadores u otros atributos (ejemplo \ref{lst:ejeCSS}). A partir de estos selectores, se definen reglas de estilo compuestas por propiedades y valores que determinan la apariencia de cada elemento.

CSS puede implementarse de distintas formas: de manera interna, dentro del mismo archivo HTML; en línea, directamente sobre los elementos; o de forma externa, mediante archivos independientes con extensión .css que se enlazan al documento HTML. Esta última forma favorece la reutilización de estilos y la organización del código en proyectos de mayor escala.

\begin{lstlisting}[language=HTML, caption={Ejemplo de código CSS.}, label={lst:ejeCSS}]
body {
  background-color: lightblue;
}
.titulo {
  color: white;
  text-align: center;
}
#paragraph {
  font-family: verdana;
  font-size: 20px;
}
\end{lstlisting}

% =======================
% |     JAVASCRIPT      |
% =======================

\subsubsection{JavaScript}

Es un lenguaje de programación utilizado en el desarrollo web para dotar de funcionalidad e interactividad a los elementos de una página. Permite implementar comportamientos dinámicos mediante la definición de eventos que se activan en respuesta a las acciones del usuario, como clics, desplazamientos o la introducción de datos \parencite[]{JSmdn}.

A través de JavaScript, es posible manipular el contenido y la estructura del documento HTML en tiempo real, así como controlar la navegación entre secciones, validar formularios, gestionar sesiones de usuario e interactuar con servicios externos mediante el intercambio de datos con servidores.

Este lenguaje se ejecuta principalmente en el navegador del cliente, aunque también puede emplearse en el lado del servidor mediante entornos de ejecución específicos. Su integración con HTML y CSS lo convierte en un componente fundamental para el desarrollo de aplicaciones web modernas.

\subsection{React.js}

React.js es una biblioteca de código abierto desarrollada por Meta Platforms, Inc. (anteriormente Facebook), utilizada para la construcción de interfaces de usuario en aplicaciones web. Se basa en el uso de componentes, los cuales son unidades independientes y reutilizables que permiten estructurar la interfaz de manera modular, facilitando el desarrollo, mantenimiento y escalabilidad de los sistemas.

Esta biblioteca emplea JavaScript y, de forma opcional, TypeScript, integrando una extensión de sintaxis conocida como JSX (o TSX en el caso de TypeScript). Esta sintaxis permite combinar estructuras similares a HTML con la lógica de programación, lo que posibilita una definición más intuitiva de la interfaz y su comportamiento en un mismo contexto.

La filosofía principal de React.js se centra en la creación de componentes reutilizables y atómicos que gestionan su propio estado (state) y propiedades (props), permitiendo que la interfaz de usuario se actualice de manera eficiente en respuesta a cambios en los datos. Este enfoque declarativo favorece el desarrollo de aplicaciones dinámicas, donde los elementos visuales pueden modificar su comportamiento e interacción según las necesidades del sistema.

\begin{figure}[ht]
    \centering
    
\includegraphics[width=0.3\linewidth]{teorico/react.png}
    \caption{Logo de React.js.}
    \label{fig:LogoReact}
\end{figure}

\subsection{Vite}

Vite es una herramienta moderna de desarrollo y construcción de aplicaciones web que facilita la ejecución, compilación y optimización del código en proyectos basados en frameworks como React. Su principal objetivo es mejorar tanto la experiencia de desarrollo como el rendimiento de las aplicaciones finales que se ejecutan en el navegador.

Vite se caracteriza por ofrecer un entorno de desarrollo rápido mediante el uso de módulos nativos del navegador, lo que permite cargar únicamente las dependencias necesarias bajo demanda. Esto contribuye a reducir los tiempos de inicio y a optimizar el consumo de recursos, evitando la carga de código innecesario. Además, durante el proceso de construcción (build), Vite genera versiones optimizadas del código, preparadas para su despliegue en entornos de producción \parencite[]{ViteJS}.

Una de sus principales características es el soporte de Hot Module Replacement (HMR), que permite reflejar en tiempo real los cambios realizados en el código sin necesidad de recargar completamente la aplicación. Esto mejora significativamente la eficiencia del desarrollo (Figura \ref{fig:Vitejs}).

\begin{figure}[ht]
    \centering
    \includegraphics[width=0.7\linewidth]{teorico/vitejs.png}
    \caption{Funcionalidad de Vitejs.}
    \label{fig:Vitejs}
\end{figure}

A diferencia de herramientas tradicionales como Babel, que se enfocan principalmente en la transformación de código, Vite integra un enfoque más completo que abarca tanto el desarrollo como la optimización del sistema, posicionándose como una alternativa moderna ampliamente adoptada en la comunidad de desarrollo web \parencite[]{BuildTools}.

\subsection{Nodejs}

Node.js es un entorno de ejecución de JavaScript, de código abierto y multiplataforma, que permite ejecutar código fuera del navegador, principalmente en el lado del servidor. Se caracteriza por su modelo de ejecución asíncrono y orientado a eventos, lo que facilita la creación de aplicaciones eficientes y altamente escalables \parencite[]{Nodejs}.

Este entorno utiliza un modelo de operaciones no bloqueantes, lo que significa que puede gestionar múltiples solicitudes de manera concurrente sin necesidad de crear múltiples hilos de ejecución. De esta forma, Node.js optimiza el uso de recursos del sistema, permaneciendo en espera cuando no existen tareas por procesar y activándose únicamente cuando se presentan eventos o solicitudes que requieren atención.

En el contexto del desarrollo de aplicaciones modernas, Node.js actúa como base para diversas herramientas y entornos, como los utilizados en proyectos con React y Vite. A través de este, es posible ejecutar servidores de desarrollo, gestionar dependencias y automatizar procesos, permitiendo simular y probar el comportamiento de las aplicaciones antes de su ejecución final en el navegador.

\begin{figure}[ht]
    \centering
    
\includegraphics[width=0.4\linewidth]{teorico/node.png}
    \caption{Logo de Node.js.}
\end{figure}